{{دیوان
| عنوان = جنگِ یرموک کا ایک واقعہ
| ماخذ = ویکی ماخذ
| ماخذ_لنک = https://ur.wikisource.org/wiki/%D8%AC%D9%86%DA%AF%D9%90_%DB%8C%D8%B1_%D9%85%D9%88%DA%A9_%DA%A9%D8%A7_%D8%A7%DB%8C%DA%A9_%D9%88%D8%A7%D9%82%D8%B9%DB%81
| مقابلہ = اقبال سائبر لائبریری (اقبال اکادمی پاکستان)
| مقابلہ_لنک = https://www.iqbalcyberlibrary.net
}}
<poem>
صف بستہ تھے عرب کے جوانانِ تیغ بند
تھی منتظر حِنا کی عروسِ زمینِ شام

اک نوجوان صُورتِ سیماب مُضطرب
آ کر ہُوا امیرِ عساکر سے ہم کلام

اے بُوعبیدہ رُخصتِ پیکار دے مجھے
لبریز ہو گیا مرے صبر و سکُوں کا جام

بے تاب ہو رہا ہوں فراقِ رسُولؐ میں
اک دم کی زندگی بھی محبّت میں ہے حرام

جاتا ہوں مَیں حضورِ رسالت پناہؐ میں
لے جاؤں گا خوشی سے اگر ہو کوئی پیام

یہ ذوق و شوق دیکھ کے پُرنم ہوئی وہ آنکھ
جس کی نگاہ تھی صفَتِ تیغِ بے نیام

بولا امیرِ فوج کہ “وہ نوجواں ہے تُو
پِیروں پہ تیرے عشق کا واجب ہے احترام

پُوری کرے خدائے محمدؐ تری مراد
کتنا بلند تیری محبّت کا ہے مقام!

پہنچے جو بارگاہِ رسُولِ امیںؐ میں تُو
کرنا یہ عرض میری طرف سے پس از سلام

ہم پر کرم کِیا ہے خدائے غیور نے
پُورے ہوئے جو وعدے کیے تھے حضورؐ نے”
</poem>
